\chapter{IoT Internetworking}

During this chapter, we will focus on the IoT internetworking aspect, which is crucial for enabling communication and data exchange between different components of an IoT system. All the different components of an IoT system need to be able to communicate with each other in order to function properly, and the IoT Things need to be able to send the measured data and receive commands and updates from the users and applications. Therefore, connectivity and communication protocols are essential for the overall functionality and performance of an IoT system.


A reader may initially think that using the Internet and the well-known TCP/IP protocol suite would be sufficient for enabling communication in an IoT system, and it is in fact a good approach as these protocols have proven to be extremely successful and widely adopted in the traditional Internet. However, there are several challenges and limitations that need to be addressed when using the Internet and TCP/IP protocols in the context of IoT systems, as an usuall Internet device is very different from an IoT device in terms of resources, capabilities, and requirements.
\begin{itemize}
    \item \ul{Data Item Size}: IoT devices often generate and exchange small data items, such as sensor readings or control commands, while the Internet and TCP/IP protocols are designed to handle larger data items, such as web pages or multimedia content.
    \item \ul{Not Optimized for Energy Efficiency}: IoT devices are often battery-powered and have limited energy resources, while the Internet and TCP/IP protocols are not optimized for energy efficiency and may require more power consumption than what is feasible for IoT devices.
    \item \ul{Large Buffer Requirements}: The Internet and TCP/IP protocols may require larger buffer sizes than what is available on IoT devices, which can lead to performance issues and increased latency.
    \item \ul{Pull vs. Push Communication}: The Internet and TCP/IP protocols are primarily designed for pull-based communication, where clients request data from servers, while IoT systems often require push-based communication, where devices need to send data to other devices or applications without waiting for a request.
\end{itemize}

In order to address these challenges and limitations, there are several IoT-specific protocols and technologies that have been developed to enable efficient and effective communication in IoT systems. There are two main approaches to connect the IoT Things to the Internet:
\begin{itemize}
    \item \ul{Application layer proxy}: 
    
    There is a proxy (usually implemented in the IoT gateway) that translates between the protocols used by the IoT Things and the protocols used by the Internet.
    \begin{itemize}
        \item The IoT Thing communicates with the proxy using dedicated IoT protocols that have nothing to do with the Internet protocols.
        \item The proxy stores the data received from the IoT Things and makes it available to the Internet using standard Internet protocols.
    \end{itemize}

    An example of this approach is when an user has to install a specific app on their smartphone in order to interact with a smart home device, such as a smart thermostat or a smart light bulb.


    \item \ul{IP Router}:
    
    In this approach, the IoT Things directly use IP, but with some optimizations and adaptations to address the challenges and limitations mentioned above. An usual router may be combined (or may replace) with the IoT gateway to enable the IoT Things to connect to the Internet using IP. That way, the IoT Things can communicate directly with other devices and applications on the Internet without the need for a proxy, which is much more scalable. There is no need for specific apps or software to interact with the IoT Things, as they can be accessed using standard Internet protocols and interfaces.

    % // TODO: Add gráfico de los protocolos?

\end{itemize}

During the rest of the course, we will explore in more detail the second approach, which is the one that is becoming more and more popular in the IoT ecosystem, as it provides better scalability and interoperability.


\section{IEEE 802.15.4}

IEEE 802.15.4 is a standard for low-rate wireless personal area networks (LR-WPANs) that is widely used in IoT applications. Its main focus is on providing low-power, low-cost and low-complexity wireless communication for IoT devices. It defines the Physical Layer (PHY) and part of the Data Link Layer (MAC). In addition, it defines multiple variants for each of the layers, which can be used in different scenarios and applications depending on the specific requirements and constraints of the IoT system.\\

Communication using this standard is based on a IEEE 802.15.4 PAN, which is a collection of devices that are connected together using the IEEE 802.15.4 standard. The members of a PAN can either be:
\begin{itemize}
    \item \ul{Full Function Devices (FFDs)}: These devices can communicate with any other device in the PAN, including other FFDs and Reduced Function Devices (RFDs). They can also take over network management activities.
    
    If two devices need to communicate with each other but they are not in range of each other, they can use an FFD to forward messages between them, which allows for multi-hop communication and extends the range of the PAN. Only FFD have this capability.
    \item \ul{Reduced Function Devices (RFDs)}: These resource-constrained devices can only communicate with one FFD in the PAN, and they cannot take over network management activities.
\end{itemize}

In addition, each of the members of a PAN can have different roles, which determine their behavior and responsibilities in the network:
\begin{itemize}
    \item \ul{PAN Coordinator}: This is the device that is responsible for managing the PAN, including tasks such as the PAN creation (by selecting a PAN ID and a channel), adding and removing devices from the PAN, and coordinating the communication between devices in the PAN. There can only be one PAN Coordinator in a PAN, and it must be an FFD.
    \item \ul{Coordinator}: This is a device that is responsible for managing a portion of the PAN, such as a cluster of devices. It can also take over some of the network management activities, such as routing and forwarding of messages. There can be multiple Coordinators in a PAN, and they must be FFDs.
    \item \ul{Members}: These are the devices that are part of the PAN and can communicate with other devices in the PAN. They can be either FFDs or RFDs, and they do not have any specific responsibilities or roles in the network.
\end{itemize}

Once the members of a PAN have been defined and assigned their roles, they can communicate with each other using one of these two types of topologies:
\begin{itemize}
    \item \ul{Star Topology}: In this topology, all the devices in the PAN communicate directly with the PAN Coordinator, which acts as a central hub for the communication. This topology is simple and easy to manage, but it can be less efficient and less scalable than other topologies, as all the communication has to go through the PAN Coordinator, which can become a bottleneck and a single point of failure.
    \item \ul{Peer-to-Peer Topology}: In this topology, the devices in the PAN can communicate directly with each other without the need for a central hub. This topology is more efficient and scalable than the star topology, as it allows for direct communication between devices, which can reduce latency and improve performance. However, it can be more complex to manage, as it requires more coordination and synchronization between devices to ensure that messages are delivered correctly and efficiently.
\end{itemize}

It should however be noted that the IEEE 802.15.4 standard does not define a protocol but only a standard for the physical and data link layers, so specific protocols as Zigbee have been developed on top of it to provide all the necessary functionalities and features for enabling communication in IoT systems.

\subsection{Physical Layer (PHY)}

The Physical Layer (PHY) of the IEEE 802.15.4 standard defines the physical characteristics and parameters of the wireless communication, such as the frequency bands where the communication can take place. It defines three different frequency bands that can be used for communication:
\begin{itemize}
    \item \ul{868 MHz Band}: This band is used in Europe and other regions, and it provides a data rate of up to 20 kbps. It has one channel available for communication.
    \item \ul{915 MHz Band}: This band is used in North America and other regions, and it provides a data rate of up to 40 kbps. It has ten channels available for communication.
    \item \ul{2.4 GHz Band}: This band is used worldwide, and it provides a data rate of up to 250 kbps. It has sixteen channels available for communication.
\end{itemize}

The services provided by the PHY layer are divided into two categories:
\begin{itemize}
    \item \ul{Data Services}: These services are used for transmitting and receiving data between devices in the PAN.
    \item \ul{Management Services}: These services are used for managing the communication and the network, such as channel scanning (to find the best channel for communication), energy detection (to measure the energy level of the channel and determine if it is busy or idle), and link quality indication (to measure the quality of the communication link between devices).
\end{itemize}

The content that is in fact transmitted over the air is the \emph{PHY Protocol Data Unit (PPDU)}, which consists of the \emph{PHY Service Data Unit (PSDU)} encapsulated with the header. The $6$ bytes Header consists of:
\begin{itemize}
    \item \ul{Synchronization Header (SHR)}, 5 bytes:
    
    This header is used for synchronization between the transmitter and the receiver, and it consists of a preamble (4 bytes) and a start of frame delimiter (1 byte).

    \item \ul{PHY Header (PHR)}, 1 byte:
    
    The first 7 bits of this header are used to indicate the length of the PSDU, while the last bit is reserved for future use. Therefore, the maximum length of the PSDU is $2^7 - 1 = 127$ bytes.
\end{itemize}

\subsection{Data Link Layer (MAC)}

\section{NB-IoT}

\section{6LowPAN}

\section{MQTT}